\section{Implementation of scattering in the $I_2$ calculation}
\label{sec:i2-scattering}

The purpose of this calculation is to reproduce Fig.~2 of
Ref.~\cite{Granelli:2021}, and then extend the same calculation by including the
$\Delta L=1$ scattering prescription used in our code.  The original paper
defines $I_2$ in its Eq.~(61),
\begin{equation}
 I_2(\kappa_1;z)
 = \int_{z_0}^{z} dz'
   \exp\!\left[-\int_{z'}^{z}d\widetilde z\,W_1(\widetilde z)\right]
   W_1(z') I_1(\kappa_1;z'),
 \label{eq:i2-definition}
\end{equation}
where $I_1$ is defined in Eq.~(59) of that reference.  The standard rates
$D_1$, $W_1$, and the equilibrium abundance $N_{N_1}^{\rm eq}$ are evaluated
with the corresponding ULYSSES functions.

\subsection{Original calculation without scattering}

Instead of evaluating the nested integrals directly, we introduce auxiliary
variables and solve the following equivalent system of first-order equations:
\begin{align}
 \frac{dN_{N_1}}{dz}
   &= -D_1\left(N_{N_1}-N_{N_1}^{\rm eq}\right),
 \label{eq:n1-no-scattering}\\
 \frac{dC}{dz}
   &= D_1\left(N_{N_1}-N_{N_1}^{\rm eq}\right)
      -\frac{W_1}{2}C,
 \label{eq:c-no-scattering}\\
 \frac{dI_1}{dz} &= C,
 \label{eq:i1-ode}\\
 \frac{dI_2}{dz} &= W_1\left(I_1-I_2\right).
 \label{eq:i2-ode}
\end{align}
Here $C$ is an auxiliary variable representing the inner, washout-damped
integral in Eq.~(59).  Equation~\eqref{eq:i2-ode} is the differential form of
Eq.~\eqref{eq:i2-definition}.  All auxiliary variables start from zero.  The
initial condition for the heavy-neutrino abundance is
\begin{equation}
 N_{N_1}(z_0)=
 \begin{cases}
  0, & \text{VIA},\\
  N_{N_1}^{\rm eq}(z_0), & \text{TIA}.
 \end{cases}
\end{equation}

\subsection{Scattering prescription}

We include the $s$- and $t$-channel scattering rates, denoted by $S_s$ and
$S_t$, following the prescription used in the repository.  Scattering enters
the calculation in two places.

First, it changes the equilibration of the heavy-neutrino abundance:
\begin{equation}
 \frac{dN_{N_1}}{dz}
 =-\left(D_1+S_s+S_t\right)
   \left(N_{N_1}-N_{N_1}^{\rm eq}\right).
 \label{eq:n1-with-scattering}
\end{equation}

Second, the inverse-decay washout is replaced by the scattering-corrected
washout
\begin{equation}
 W_{\rm sct}
 =W_1\left[
 1+\frac{1}{D_1}
 \left(
  2\frac{N_{N_1}}{N_{N_1}^{\rm eq}}S_s+4S_t
 \right)
 \right].
 \label{eq:scattering-washout}
\end{equation}

The decay rate $D_1$, rather than $D_1+S_s+S_t$, is retained in the source
term.  Thus, scattering changes the production of $N_1$ and the washout, but
does not act as an additional CP-violating source.  The system solved with
scattering is therefore
\begin{align}
 \frac{dN_{N_1}}{dz}
   &= -\left(D_1+S_s+S_t\right)
      \left(N_{N_1}-N_{N_1}^{\rm eq}\right),\\
 \frac{dC}{dz}
   &= D_1\left(N_{N_1}-N_{N_1}^{\rm eq}\right)
      -\frac{W_{\rm sct}}{2}C,\\
 \frac{dI_1}{dz} &= C,\\
 \frac{dI_2}{dz} &= W_{\rm sct}\left(I_1-I_2\right).
 \label{eq:i2-system-scattering}
\end{align}
Consequently, the scattering extension amounts to
\begin{equation}
 D_1\rightarrow D_1+S_s+S_t
 \quad\text{only in the $N_1$ equation},
 \qquad
 W_1\rightarrow W_{\rm sct}
 \quad\text{in all washout factors}.
\end{equation}
The $D_1(N_{N_1}-N_{N_1}^{\rm eq})$ source is unchanged.

\subsection{Numerical implementation}

The calculation is implemented in \texttt{src/benchmark\_analysis.py} in the
functions \texttt{scan\_i2} and \texttt{\_i2\_at\_kappa}.  For every value of
$\kappa_1$, the four coupled equations are integrated with
\texttt{scipy.integrate.solve\_ivp} using the LSODA method from
$z_0=10^{-3}$ to $z_f=10^3$.  The ULYSSES methods \texttt{D1}, \texttt{W1},
and \texttt{N1Eq} supply the standard rates and equilibrium abundance.  The
functions \texttt{scat\_Ss}, \texttt{scat\_St}, and
\texttt{washout\_sct}, defined in \texttt{src/scattering.py}, supply the
scattering rates and Eq.~\eqref{eq:scattering-washout}.

The rates are linear in $\kappa_1$.  They are evaluated once at
$\kappa_1=1$ on a logarithmic $z$ grid and rescaled for each point of the
$\kappa_1$ scan.  The same numerical setup and initial conditions are used in
the calculations with and without scattering, so the comparison isolates the
effect of the scattering prescription.

For clarity, this is not a built-in ULYSSES evaluation of $I_2$.  ULYSSES is
used to obtain the standard ingredients $D_1$, $W_1$, and
$N_{N_1}^{\rm eq}$; the coupled system defining $I_1$ and $I_2$, together with
its scattering extension, is implemented separately in this repository.

